% ══════════════════════════════════════════════════════════════
%  preamble.tex — Rust Blockchain: A Full-Stack Implementation Guide
%  Professional typesetting configuration for Crown Quarto
% ══════════════════════════════════════════════════════════════

% ── Overfull tolerance ──────────────────────────────────────
%  Suppress overfull hbox warnings under 3pt (cosmetically invisible).
%  Code listings with long lines will still overflow slightly into the margin.
\hfuzz=3pt

% ── Encoding & language ──────────────────────────────────────
\usepackage[T1]{fontenc}
\usepackage[utf8]{inputenc}
\usepackage[english]{babel}

% ── Fonts ────────────────────────────────────────────────────
%  Body: TeX Gyre Pagella (Palatino family — elegant, highly readable)
%  Sans:  TeX Gyre Heros (Helvetica family — clean headings)
%  Mono:  TeX Gyre Cursor scaled to match (Courier family, for inline code)
%  Math:  Matching Palatino math via mathpazo
\usepackage{tgpagella}           % Serif body
\usepackage{tgheros}             % Sans-serif for headings
\usepackage{tgcursor}               % Monospace for code
\usepackage[sc]{mathpazo}        % Palatino math + small caps
\linespread{1.05}                % Pagella reads better with slightly looser leading

% ── Drop caps ──────────────────────────────────────────────
\usepackage{lettrine}
\renewcommand{\LettrineFontHook}{\sffamily\bfseries\color{brand}}
\setcounter{DefaultLines}{3}
\renewcommand{\DefaultLoversize}{0.1}
\renewcommand{\DefaultLhang}{0.0}
\setlength{\DefaultFindent}{3pt}
\setlength{\DefaultNindent}{0pt}

% ── Microtypography ──────────────────────────────────────────
\usepackage[
  activate={true,nocompatibility},
  final,
  tracking=true,
  kerning=true,
  spacing=true,
  factor=1100,
  stretch=10,
  shrink=10
]{microtype}
\microtypecontext{spacing=nonfrench}
\frenchspacing                   % Single space after periods (modern style)

% ── Page geometry (Crown Quarto: 7.44" × 9.69") ─────────────
%  IngramSpark print-ready: 0.125" bleed on all sides
%  Stock = trim + 2×bleed = 7.69" × 9.94"
\setstocksize{9.94in}{7.69in}
\settrimmedsize{9.69in}{7.44in}{*}
\settrims{0.125in}{0.125in}         % 0.125" bleed offset (top, foredge)

% Generous margins for readability; centered text block on every page
\settypeblocksize{7.5in}{5.0in}{*}
\setbinding{0pt}                % No binding shift — text block is perfectly centered
\setlrmargins{*}{*}{1}          % Equal inner/outer margins
\setulmargins{0.85in}{*}{*}     % 0.85" top; bottom computed
\setheadfoot{\baselineskip}{2\baselineskip}
\setheaderspaces{*}{1.5\baselineskip}{*}
\checkandfixthelayout

% ── Orphan heading prevention ────────────────────────────────
%  Ensure section/subsection headings never land at the page bottom
%  without enough following content on the same page.
\usepackage{needspace}
\usepackage{etoolbox}
\pretocmd{\section}{\needspace{1.8in}}{}{}
\pretocmd{\subsection}{\needspace{1.3in}}{}{}
\pretocmd{\subsubsection}{\needspace{1in}}{}{}

% Prevent orphan chapter entries at the bottom of TOC/LOL pages:
% Add vertical space before chapter-level TOC entries that also acts
% as a needspace trigger — if less than 4 lines remain, push to next page.
\setlength{\cftbeforechapterskip}{1.2em}          % visual gap before each chapter in TOC
\renewcommand{\cftchapterbreak}{\needspace{6\baselineskip}}  % prevent orphan TOC entries

% ── Part entries in TOC: large, bold, red, with generous spacing ──
\renewcommand{\cftpartfont}{\sffamily\LARGE\bfseries\color{brand}}
\renewcommand{\cftpartpagefont}{\sffamily\LARGE\bfseries\color{brand}}
\renewcommand{\cftpartleader}{\hfill}    % clean fill (no dots) before page number
\setlength{\cftbeforepartskip}{28pt}     % generous space above Part in TOC
% Force a page break before each Part entry in the TOC to prevent orphans
% and give Part headings visual prominence
\renewcommand{\cftpartbreak}{\clearpage}

% Discourage widows (last line of paragraph alone at top of page)
% and orphans (first line of paragraph alone at bottom of page)
\widowpenalty=10000
\clubpenalty=10000

% Prevent isolated lines before/after display environments (code, figures)
\predisplaypenalty=10000
\postdisplaypenalty=500

% ── Colors ───────────────────────────────────────────────────
\usepackage[dvipsnames,svgnames,x11names,table]{xcolor}
% ──────────────────────────────────────────────────────────────
%  colors.tex — Unified color palette for Rust Blockchain book
% ──────────────────────────────────────────────────────────────
%  All colors are defined here.  Every other file references
%  these names so the palette can be changed in one place.
% ──────────────────────────────────────────────────────────────

% ── Brand / accent ───────────────────────────────────────────
\definecolor{brand}{HTML}{B7410E}        % Rust-orange (chapter numbers, rules)
\definecolor{brandlight}{HTML}{F4E0D0}   % Tinted background for brand accents
\definecolor{brandgray}{HTML}{3B3B3B}    % Near-black for body text

% ── Structural ───────────────────────────────────────────────
\definecolor{chapternum}{HTML}{B7410E}   % Large chapter number
\definecolor{chaptertitle}{HTML}{1A1A1A} % Chapter title text
\definecolor{sectioncolor}{HTML}{2C5F7C} % Section headings (dark teal)
\definecolor{rulecolor}{HTML}{B7410E}    % Decorative horizontal rules

% ── Code blocks ──────────────────────────────────────────────
\definecolor{codebg}{HTML}{F7F7F7}       % Code background (warm near-white)
\definecolor{codeframe}{HTML}{DDDDDD}    % Code block border
\definecolor{codekw}{HTML}{7C2D92}       % Keywords (purple)
\definecolor{codecomment}{HTML}{6A737D}  % Comments (muted gray)
\definecolor{codestring}{HTML}{22863A}   % Strings (green)
\definecolor{codenumber}{HTML}{005CC5}   % Numbers and constants (blue)
\definecolor{codelinenum}{HTML}{AAAAAA}  % Line numbers

% ── Callout boxes ────────────────────────────────────────────
% Info family (blue-gray): Methods involved, Source, Scope
\definecolor{infobg}{HTML}{EBF5FB}
\definecolor{infoframe}{HTML}{3498DB}
\definecolor{infotitle}{HTML}{1A5276}

% Note family (blue): Note
\definecolor{notebg}{HTML}{E8EAF6}
\definecolor{noteframe}{HTML}{5C6BC0}
\definecolor{notetitle}{HTML}{283593}

% Tip family (green): Tip
\definecolor{tipbg}{HTML}{E8F5E9}
\definecolor{tipframe}{HTML}{4CAF50}
\definecolor{tiptitle}{HTML}{1B5E20}

% Learning family (green): What you will learn
\definecolor{learnbg}{HTML}{E0F2F1}
\definecolor{learnframe}{HTML}{26A69A}
\definecolor{learntitle}{HTML}{004D40}

% Warning family (amber): Warning, Troubleshooting
\definecolor{warnbg}{HTML}{FFF8E1}
\definecolor{warnframe}{HTML}{FFA000}
\definecolor{warntitle}{HTML}{E65100}

% Important family (red): Important
\definecolor{importantbg}{HTML}{FCE4EC}
\definecolor{importantframe}{HTML}{E53935}
\definecolor{importanttitle}{HTML}{B71C1C}

% Prereq family (orange): Prerequisites
\definecolor{prereqbg}{HTML}{FFF3E0}
\definecolor{prereqframe}{HTML}{FB8C00}
\definecolor{prereqtitle}{HTML}{E65100}

% Example family (teal): See it in action, Implementation example, Framework Comparison
\definecolor{examplebg}{HTML}{E0F7FA}
\definecolor{exampleframe}{HTML}{00ACC1}
\definecolor{exampletitle}{HTML}{006064}

% Reference family (gray): See the full implementation, Companion Chapter
\definecolor{refbg}{HTML}{ECEFF1}
\definecolor{refframe}{HTML}{78909C}
\definecolor{reftitle}{HTML}{37474F}

% Checkpoint family (purple): Checkpoint
\definecolor{checkpointbg}{HTML}{F3E5F5}
\definecolor{checkpointframe}{HTML}{AB47BC}
\definecolor{checkpointtitle}{HTML}{6A1B9A}


% ── Graphics and floats ──────────────────────────────────────
\usepackage{graphicx}
\usepackage{float}
\usepackage{caption}
\captionsetup{
  font={small,sf},
  labelfont={bf,color=sectioncolor},
  format=hang,
  margin=10pt,
}
\captionsetup[listing]{
  position=above,
  skip=4pt,
}

% ── Tables ───────────────────────────────────────────────────
\usepackage{booktabs}            % Professional table rules
\usepackage{longtable}           % Multi-page tables
\usepackage{array}               % Enhanced column types
\usepackage{tabularx}            % Flexible-width columns
\renewcommand{\arraystretch}{1.3}% Slightly more row spacing

% ── Code listings (minted) ───────────────────────────────────
% minted v3+ (TeX Live 2024+) auto-detects the output directory;
% the old outputdir=build option is no longer accepted.
\usepackage{minted}
\usemintedstyle{friendly}        % Clean, readable style

% Default minted settings for all languages
\setminted{
  fontsize=\small,
  bgcolor=codebg,
  frame=leftline,
  framerule=2pt,
  rulecolor=codeframe,
  breaklines=true,
  breakanywhere=true,
  breaksymbol={},
  breakbeforesymbolpre={},
  breakaftersymbolpre={},
  tabsize=4,
  numbersep=8pt,
  xleftmargin=14pt,
}

% Rust-specific overrides
\setminted[rust]{
  linenos=true,
  escapeinside=@@,             % Allow LaTeX inside @...@ in code
}

% ── Inline code styling (background pill) ─────────────────────
%  Renders inline code like  \code{BlockHeader}  with a subtle
%  gray pill background so it pops from prose at a glance.
%  Uses \colorbox which is robust in captions and moving arguments.
\definecolor{inlinebg}{HTML}{F0F0F0}       % Subtle warm gray
\definecolor{INLINEBG}{HTML}{F0F0F0}       % Alias for hyperref bookmark context
\newcommand{\code}[1]{%
  \texorpdfstring
    {\colorbox{inlinebg}{\texttt{\small\vphantom{Ag}#1}}}%
    {#1}%
}

% ── Listing float environment ─────────────────────────────────
\newfloat{listing}{htbp}{lol}[chapter]
\floatname{listing}{Listing}
\renewcommand{\listoflistings}{\listof{listing}{List of Listings}}

% Better spacing around listings
\setlength{\intextsep}{12pt plus 4pt minus 2pt}   % Space around in-text floats
\setlength{\textfloatsep}{16pt plus 4pt minus 3pt} % Space around top/bottom floats

% ── Callout boxes (tcolorbox) ────────────────────────────────
\usepackage[most]{tcolorbox}
% ──────────────────────────────────────────────────────────────
%  callouts.tex — tcolorbox styles for all 16 canonical callout types
% ──────────────────────────────────────────────────────────────
%  Usage:  \begin{notebox} ... \end{notebox}
%          \begin{infobox}{Methods involved} ... \end{infobox}
%  Each box has a visual icon prefix in its title bar for instant
%  recognition while scanning or flipping pages.
% ──────────────────────────────────────────────────────────────

\tcbuselibrary{skins,breakable}

% ── Icon definitions (using text symbols for max compatibility) ─
\newcommand{\iconinfo}{\textsf{\textbf{i}}\kern4pt}
\newcommand{\iconnote}{\raisebox{0.5pt}{\footnotesize$\triangleright$}\kern4pt}
\newcommand{\icontip}{\raisebox{-0.5pt}{\footnotesize$\star$}\kern4pt}
\newcommand{\iconlearn}{\raisebox{-0.5pt}{\footnotesize$\vartriangleright$}\kern4pt}
\newcommand{\iconwarn}{\raisebox{0pt}{\footnotesize$\blacktriangle$}\kern4pt}
\newcommand{\iconimportant}{\raisebox{0pt}{\footnotesize$\blacklozenge$}\kern4pt}
\newcommand{\iconprereq}{\raisebox{-0.5pt}{\footnotesize$\bullet\!\bullet\!\bullet$}\kern4pt}
\newcommand{\iconexample}{\raisebox{-0.5pt}{\footnotesize$\rhd$}\kern4pt}
\newcommand{\iconref}{\raisebox{0pt}{\footnotesize$\hookrightarrow$}\kern4pt}
\newcommand{\iconcheckpoint}{\raisebox{0pt}{\footnotesize$\checkmark$}\kern4pt}
\newcommand{\icontrouble}{\raisebox{0pt}{\footnotesize$\wr$}\kern4pt}

% ── Shared base style ────────────────────────────────────────
\tcbset{
  calloutbase/.style={
    enhanced,
    breakable,
    left=10pt,
    right=10pt,
    top=8pt,
    bottom=8pt,
    boxrule=0pt,
    leftrule=3pt,
    arc=2pt,
    fonttitle=\sffamily\bfseries\small,
    fontupper=\small,
    before upper={\parindent=0pt},
    before skip=10pt plus 2pt,
    after skip=10pt plus 2pt,
  },
}

% ── Info (Methods involved, Source, Scope) ───────────────────
\newtcolorbox{infobox}[1][]{
  calloutbase,
  colback=infobg,
  colframe=infoframe,
  coltitle=infotitle,
  title={\iconinfo #1},
}

% ── Note ─────────────────────────────────────────────────────
\newtcolorbox{notebox}[1][Note]{
  calloutbase,
  colback=notebg,
  colframe=noteframe,
  coltitle=notetitle,
  title={\iconnote #1},
}

% ── Tip ──────────────────────────────────────────────────────
\newtcolorbox{tipbox}[1][Tip]{
  calloutbase,
  colback=tipbg,
  colframe=tipframe,
  coltitle=tiptitle,
  title={\icontip #1},
}

% ── Learning objectives (What you will learn) ────────────────
\newtcolorbox{learnbox}[1][What you will learn in this chapter]{
  calloutbase,
  colback=learnbg,
  colframe=learnframe,
  coltitle=learntitle,
  title={\iconlearn #1},
}

% ── Warning ──────────────────────────────────────────────────
\newtcolorbox{warnbox}[1][Warning]{
  calloutbase,
  colback=warnbg,
  colframe=warnframe,
  coltitle=warntitle,
  title={\iconwarn #1},
}

% ── Important ────────────────────────────────────────────────
\newtcolorbox{importantbox}[1][Important]{
  calloutbase,
  colback=importantbg,
  colframe=importantframe,
  coltitle=importanttitle,
  title={\iconimportant #1},
}

% ── Prerequisites ────────────────────────────────────────────
\newtcolorbox{prereqbox}[1][Prerequisites]{
  calloutbase,
  colback=prereqbg,
  colframe=prereqframe,
  coltitle=prereqtitle,
  title={\iconprereq #1},
}

% ── Example (See it in action, Implementation example, Framework Comparison) ─
\newtcolorbox{examplebox}[1][]{
  calloutbase,
  colback=examplebg,
  colframe=exampleframe,
  coltitle=exampletitle,
  title={\iconexample #1},
}

% ── Reference (See the full implementation, Companion Chapter) ─
\newtcolorbox{refbox}[1][]{
  calloutbase,
  colback=refbg,
  colframe=refframe,
  coltitle=reftitle,
  title={\iconref #1},
}

% ── Checkpoint ───────────────────────────────────────────────
\newtcolorbox{checkpointbox}[1][Checkpoint]{
  calloutbase,
  colback=checkpointbg,
  colframe=checkpointframe,
  coltitle=checkpointtitle,
  title={\iconcheckpoint #1},
}

% ── Troubleshooting (uses warning palette) ───────────────────
\newtcolorbox{troublebox}[1][Troubleshooting]{
  calloutbase,
  colback=warnbg,
  colframe=warnframe,
  coltitle=warntitle,
  title={\icontrouble #1},
}


% ── Cross-references & hyperlinks ────────────────────────────
\usepackage{hyperref}
\hypersetup{
  colorlinks=true,
  linkcolor=sectioncolor,
  citecolor=sectioncolor,
  urlcolor=brand,
  bookmarksnumbered=true,
  bookmarksopen=true,
  bookmarksopenlevel=1,
  pdfstartview=FitH,
  pdftitle={Rust Blockchain: A Full-Stack Implementation Guide},
  pdfauthor={Bill Kunyiha},
  pdfsubject={Blockchain development with Rust},
  pdfkeywords={Rust, Blockchain, Bitcoin, Tokio, Axum, Iced, Tauri},
}
\usepackage[nameinlink,capitalize]{cleveref} % Smart cross-references (\cref)
\crefname{listing}{Listing}{Listings}
\Crefname{listing}{Listing}{Listings}
\crefformat{chapter}{Chapter~#2#1#3}
\crefformat{section}{Section~#2#1#3}
\crefformat{listing}{Listing~#2#1#3}

% ── Bibliography ─────────────────────────────────────────────
% Uncomment when bibliography.bib is ready:
% \usepackage[backend=biber,style=numeric,sorting=none]{biblatex}
% \addbibresource{backmatter/bibliography.bib}

% ── Index ────────────────────────────────────────────────────
\usepackage{makeidx}
\makeindex

% ── QR code ──────────────────────────────────────────────────
\usepackage{qrcode}

% ── Math ─────────────────────────────────────────────────────
\usepackage{amsmath}
\usepackage{amssymb}

% ── Epigraphs (chapter-opening quotes) ───────────────────────
\usepackage{epigraph}
\setlength{\epigraphwidth}{0.75\textwidth}
\setlength{\epigraphrule}{0pt}             % No rule under epigraph
\renewcommand{\epigraphflush}{flushright}
\renewcommand{\sourceflush}{flushright}
\renewcommand{\epigraphsize}{\small\itshape}

% ── Chapter mini-TOC ──────────────────────────────────────────
%  A manual shaded box listing a chapter's sections for navigation.
%  Usage:  \chapteroutline{
%    \item Section title one
%    \item Section title two
%  }
\newenvironment{chapteroutline}{%
  \begin{tcolorbox}[
    enhanced, breakable,
    colback=codebg, colframe=codeframe,
    boxrule=0.5pt, arc=3pt,
    left=12pt, right=12pt, top=8pt, bottom=8pt,
    title={\sffamily\small\bfseries\color{sectioncolor} In this chapter},
    fonttitle=\sffamily\small\bfseries,
    coltitle=sectioncolor,
    before skip=8pt, after skip=14pt,
  ]%
  \small\sffamily
  \begin{enumerate}[label={\color{brand}\arabic*.},leftmargin=1.5em,nosep,itemsep=2pt]%
}{%
  \end{enumerate}%
  \end{tcolorbox}%
}

% ══════════════════════════════════════════════════════════════
%  MEMOIR CHAPTER & SECTION STYLES
% ══════════════════════════════════════════════════════════════

% ── Chapter style: custom "blockchain" ───────────────────────
\makechapterstyle{blockchain}{%
  \renewcommand{\chapnamefont}{\sffamily\Large\bfseries\color{sectioncolor}}%
  \renewcommand{\chapnumfont}{\sffamily\fontsize{56}{60}\selectfont\color{chapternum}}%
  \renewcommand{\chaptitlefont}{\sffamily\HUGE\bfseries\color{chaptertitle}}%
  \renewcommand{\printchaptername}{}%
  \renewcommand{\chapternamenum}{}%
  \renewcommand{\afterchapternum}{}%
  \renewcommand{\printchapternum}{%
    \begin{adjustwidth}{}{-10pt}%
      \hfill\chapnumfont\thechapter%
    \end{adjustwidth}%
    \vspace{-18pt}%
    \hfill{\color{rulecolor}\rule{0.85\textwidth}{2pt}}%
    \vspace{12pt}%
  }%
  \renewcommand{\printchaptertitle}[1]{%
    \begin{adjustwidth}{}{0pt}%
      \raggedright\chaptitlefont ##1%
    \end{adjustwidth}%
  }%
  \renewcommand{\afterchaptertitle}{%
    \vskip 20pt%
  }%
}
\chapterstyle{blockchain}

% ── Section styles ───────────────────────────────────────────
\setsecheadstyle{\sffamily\Large\bfseries\color{sectioncolor}}
\setsubsecheadstyle{\sffamily\large\bfseries\color{sectioncolor}}
\setsubsubsecheadstyle{\sffamily\normalsize\bfseries\color{sectioncolor}}

% ── Paragraph indentation ────────────────────────────────────
\setlength{\parindent}{1.2em}    % Standard paragraph indent
\setlength{\parskip}{0pt plus 1pt} % No skip between paragraphs (book style)

% ── Headers and footers ─────────────────────────────────────
%  Verso (even): page number left, chapter title right
%  Recto (odd):  section title left, page number right
%  Thin rule underneath for clean separation
\makepagestyle{blockchain}
\makeevenhead{blockchain}%
  {\small\sffamily\thepage}%
  {}%
  {\small\sffamily\itshape\textcolor{sectioncolor}{\leftmark}}
\makeoddhead{blockchain}%
  {\small\sffamily\itshape\textcolor{sectioncolor}{\rightmark}}%
  {}%
  {\small\sffamily\thepage}
\makeevenfoot{blockchain}{}{}{}
\makeoddfoot{blockchain}{}{}{}
\makeheadrule{blockchain}{\textwidth}{0.3pt}
\pagestyle{blockchain}

% Alias for plain pages (chapter openings)
\copypagestyle{plain}{blockchain}
\makeevenhead{plain}{}{}{}
\makeoddhead{plain}{}{}{}
\makeevenfoot{plain}{}{\small\sffamily\thepage}{}
\makeoddfoot{plain}{}{\small\sffamily\thepage}{}

% ── Part page style ──────────────────────────────────────────
\renewcommand{\partnamefont}{\sffamily\Huge\bfseries\color{sectioncolor}}
\renewcommand{\partnumfont}{\sffamily\fontsize{72}{76}\selectfont\color{chapternum}}
\renewcommand{\parttitlefont}{\sffamily\HUGE\bfseries\color{chaptertitle}}

% Part overview: add description text below part title on the SAME page.
% Usage:  \partoverview{Chapters 1--8 introduce...}
%  Place immediately after \part{...}
\makeatletter
\newcommand{\bk@parttext}{}
\newcommand{\partoverview}[1]{%
  \gdef\bk@parttext{#1}%
}
\renewcommand{\afterpartskip}{%
  \ifx\bk@parttext\empty\else
    \vspace{20pt}%
    \begin{center}
      \begin{minipage}{0.75\textwidth}
        \centering\sffamily\large\color{sectioncolor!80}%
        \bk@parttext
      \end{minipage}
    \end{center}%
    \gdef\bk@parttext{}%
  \fi
  \vfil\newpage
}
\makeatother

% ── Miscellaneous ────────────────────────────────────────────
\usepackage{changepage}          % adjustwidth environment
\usepackage{enumitem}            % List customization
\setlist{nosep,leftmargin=*}     % Tight lists, no extra spacing
\setlist[itemize]{label={\color{brand}\textbullet}}
\setlist[enumerate]{label={\color{brand}\arabic*.}}

% ── Widow, orphan, and page-break control ─────────────────────
\widowpenalty=10000
\clubpenalty=10000
\brokenpenalty=4991               % Discourage hyphenation across pages
\predisplaypenalty=10000          % Avoid page breaks just before display math
\postdisplaypenalty=0

% ── Vertical rhythm helpers ───────────────────────────────────
%  Consistent spacing before/after section headings
\setbeforesecskip{-2.5ex plus -1ex minus -.2ex}
\setaftersecskip{1.3ex plus .2ex}
\setbeforesubsecskip{-2.25ex plus -1ex minus -.2ex}
\setaftersubsecskip{1ex plus .2ex}

% ── "Key term" command for first-use definitions ──────────────
%  Prints bold + adds to index:  \keyterm{UTXO} → UTXO (bold, indexed)
\newcommand{\keyterm}[1]{\textbf{#1}\index{#1}}

% ── Separator between logical sections within a chapter ───────
%  A centered *** ornament for thematic breaks
\newcommand{\scenbreak}{\bigskip\par\noindent\centerline{*\quad*\quad*}\bigskip\par}

% ── Margin notes ────────────────────────────────────────────
%  Usage:  \marginnote{See also §4.2}  or  \marginnote{New in v0.3}
%  Prints a small annotation in the outer margin.
\newcommand{\marginnote}[1]{%
  \marginpar{\raggedright\footnotesize\sffamily\color{sectioncolor}#1}%
}

% ── Code annotation callouts (Manning-style) ────────────────
%  Numbered circles in code listings with matching explanations.
%  Usage in code: // \cmark{1}
%  Usage below listing:
%    \begin{codeannotations}
%      \annotate{1}{Creates the block header with timestamp}
%      \annotate{2}{Runs proof-of-work to find valid nonce}
%    \end{codeannotations}
\newcommand{\cmark}[1]{%
  \tikz[baseline=(char.base)]{%
    \node[circle,fill=brand,text=white,inner sep=1pt,font=\sffamily\bfseries\tiny,minimum size=12pt] (char) {#1};%
  }%
}
\newenvironment{codeannotations}{%
  \vspace{-4pt}%
  \begin{description}[
    font=\normalfont,
    leftmargin=2.5em,
    labelwidth=2em,
    labelsep=0.5em,
    nosep,
    itemsep=3pt,
  ]%
}{%
  \end{description}%
  \vspace{4pt}%
}
\newcommand{\annotate}[2]{%
  \item[\cmark{#1}] {\small #2}%
}

% ── Muted code line numbers ─────────────────────────────────
%  Override minted line number color to muted gray
\renewcommand{\theFancyVerbLine}{%
  {\small\sffamily\textcolor{codelinenum}{\arabic{FancyVerbLine}}}%
}

% ── Chapter-end summary box ─────────────────────────────────
%  Usage:
%    \begin{chaptersummary}
%      \item Blocks link via \code{pre\_block\_hash}
%      \item The UTXO set tracks spendability
%    \end{chaptersummary}
\newtcolorbox{chaptersummarybox}{
  enhanced, breakable,
  colback=brandlight,
  colframe=brand,
  boxrule=1pt,
  arc=4pt,
  left=12pt, right=12pt, top=8pt, bottom=8pt,
  title={\sffamily\bfseries\color{white} Key Takeaways},
  fonttitle=\sffamily\bfseries,
  coltitle=white,
  colbacktitle=brand,
  before skip=14pt, after skip=10pt,
}
\newenvironment{chaptersummary}{%
  \begin{chaptersummarybox}%
  \small
  \begin{itemize}[nosep,itemsep=3pt,leftmargin=1.2em,label={\color{brand}\checkmark}]%
}{%
  \end{itemize}%
  \end{chaptersummarybox}%
}

% ── Glossary support ────────────────────────────────────────
%  Using glossaries-extra for auto-linked term definitions.
%  Terms are defined in backmatter/glossary.tex and used via
%  \gls{utxo}, \glspl{utxo} (plural), \Gls{utxo} (capitalized)
\usepackage[toc,acronym,nonumberlist,nogroupskip]{glossaries-extra}
\setabbreviationstyle[acronym]{long-short}
\renewcommand{\glsnamefont}[1]{\sffamily\bfseries\color{sectioncolor}#1}
\makeglossaries

% ── Thumb index tabs ────────────────────────────────────────
%  Colored edge tabs by part, appearing on recto page trim edges.
%  Each \setparttab{color} call updates the active tab color.
\usepackage{tikz}
\usetikzlibrary{calc,positioning,shapes.geometric,arrows.meta,backgrounds,fit}

\newcounter{thumbtab}
\newlength{\thumbtabheight}
\setlength{\thumbtabheight}{1.8cm}
\newlength{\thumbtabwidth}
\setlength{\thumbtabwidth}{0.5cm}
\newlength{\thumbtabsep}
\setlength{\thumbtabsep}{0.1cm}

% Color for current part's tab — updated by \setparttab
\colorlet{currenttab}{brand}

\newcommand{\setparttab}[2]{%
  \setcounter{thumbtab}{#1}%
  \colorlet{currenttab}{#2}%
}

% Draw the thumb tab on each page via eso-pic
\usepackage{eso-pic}
\newcommand{\thumbmark}{%
  \ifodd\value{page}%
    \begin{tikzpicture}[remember picture,overlay]
      \fill[currenttab,rounded corners=2pt]
        ([yshift=-{(\value{thumbtab}-1)*(\thumbtabheight+\thumbtabsep)+0.5cm}]current page.north east)
        rectangle
        +(-\thumbtabwidth,-\thumbtabheight);
    \end{tikzpicture}%
  \fi
}
\AddToShipoutPictureBG{\thumbmark}

% ── Alternating table row shading ───────────────────────────
%  Usage:  \rowcolors{2}{white}{tableshade} before a tabular
\definecolor{tableshade}{HTML}{F5F6FA}  % Very subtle blue-gray

% ── TikZ diagram style guide ───────────────────────────────
%  Consistent styles for all architecture diagrams in the book.
\tikzset{
  % Node styles
  blockchain node/.style={
    rectangle, rounded corners=4pt,
    draw=sectioncolor, fill=codebg,
    thick, font=\sffamily\small,
    minimum height=2em, minimum width=4em,
    text centered,
  },
  blockchain component/.style={
    rectangle, rounded corners=6pt,
    draw=brand, fill=brandlight,
    thick, font=\sffamily\small\bfseries,
    minimum height=2.5em, minimum width=5em,
    text centered,
  },
  blockchain data/.style={
    rectangle,
    draw=codeframe, fill=codebg,
    font=\ttfamily\small,
    minimum height=2em,
    inner sep=6pt,
  },
  blockchain group/.style={
    rectangle, rounded corners=8pt,
    draw=sectioncolor!50, dashed,
    fill=sectioncolor!5,
    inner sep=10pt,
    label={[font=\sffamily\small\bfseries\color{sectioncolor}]above:#1},
  },
  % Arrow styles
  blockchain arrow/.style={
    -{Stealth[length=6pt,width=4pt]},
    thick, color=sectioncolor,
  },
  blockchain dotted/.style={
    -{Stealth[length=6pt,width=4pt]},
    thick, dashed, color=sectioncolor!60,
  },
  data flow/.style={
    -{Stealth[length=6pt,width=4pt]},
    thick, color=brand,
  },
}

% ── Pull quotes ─────────────────────────────────────────────
%  Large, styled callout for key one-liner insights mid-chapter.
%  Usage:  \pullquote{The UTXO set is derived state, not history.}
\newcommand{\pullquote}[1]{%
  \begin{center}
  \begin{minipage}{0.82\textwidth}
    \vspace{8pt}%
    {\color{brand}\rule{\linewidth}{1.5pt}}\\[6pt]
    {\large\sffamily\itshape\color{sectioncolor}%
      \hspace{0.5em}``#1''}\\[6pt]
    {\color{brand}\rule{\linewidth}{1.5pt}}%
    \vspace{8pt}%
  \end{minipage}
  \end{center}
}

% ── Deep dive sidebar ───────────────────────────────────────
%  A full-width shaded aside for optional advanced content.
%  Readers can skip without losing the main thread.
%  Usage:
%    \begin{deepdive}{Why Merkle trees matter}
%      Advanced explanation here...
%    \end{deepdive}
\newtcolorbox{deepdive}[1]{
  enhanced, breakable,
  colback=sectioncolor!5,
  colframe=sectioncolor!40,
  boxrule=0.5pt,
  arc=4pt,
  left=12pt, right=12pt, top=8pt, bottom=8pt,
  title={\sffamily\bfseries\small\color{white}\raisebox{-0.5pt}{\footnotesize$\vartriangleright$}\kern4pt Deep Dive:\enspace #1},
  fonttitle=\sffamily\bfseries\small,
  coltitle=white,
  colbacktitle=sectioncolor,
  fontupper=\small,
  before skip=12pt, after skip=12pt,
  borderline south={1pt}{0pt}{sectioncolor!40},
}

% ── Step-by-step tutorial environment ───────────────────────
%  Distinct visual style for tutorial walkthroughs.
%  Usage:
%    \begin{steps}
%      \step{Clone the repository}
%        Run \code{git clone ...} to get started.
%      \step{Build the project}
%        Run \code{cargo build --release}.
%    \end{steps}
\newcounter{stepcount}
\newsavebox{\stepbox}
\newcommand{\makesteplabel}[1]{%
  \sbox{\stepbox}{\colorbox{brand}{\color{white}\sffamily\bfseries\small\strut\enspace Step #1\enspace}}%
  \usebox{\stepbox}%
}
\newenvironment{steps}{%
  \setcounter{stepcount}{0}%
  \begin{list}{}{%
    \setlength{\leftmargin}{4.5em}%
    \setlength{\labelwidth}{4.2em}%
    \setlength{\labelsep}{0.3em}%
    \setlength{\itemsep}{8pt plus 2pt}%
    \setlength{\parsep}{3pt plus 1pt}%
    \setlength{\topsep}{6pt plus 2pt}%
  }%
}{%
  \end{list}%
}
\newcommand{\step}[1]{%
  \stepcounter{stepcount}%
  \item[\makesteplabel{\thestepcount}]%
  {\sffamily\bfseries\color{sectioncolor}#1}\par\nopagebreak
}

% ── Footnote styling ────────────────────────────────────────
%  Brand-colored markers with hanging indent and thin separator
\renewcommand{\footnoterule}{%
  \kern-3pt
  {\color{brand}\hrule width 0.3\textwidth height 0.4pt}%
  \kern 2.6pt
}
\renewcommand{\thefootnote}{%
  {\color{brand}\sffamily\arabic{footnote}}%
}
\renewcommand{\footnotesize}{\small}

% ── API reference table style ──────────────────────────────
%  Standardized format for method signatures.
%  Usage:  \apitable{StructName}{rows}  where rows use \apirow
%  Example:
%    \apitable{BlockchainService}{%
%      \apirow{add\_block}{block: Block}{Result<()>}{Validates and appends a block}
%      \apirow{get\_tip}{---}{Option<Block>}{Returns the chain tip}
%    }
\newcommand{\apitable}[2]{%
  \begin{table}[H]
  \caption{API Reference: \code{#1}}
  \centering\small
  \rowcolors{2}{white}{tableshade}%
  \begin{tabular}{@{} l l l p{4.8cm} @{}}
  \toprule
  \textbf{Method} & \textbf{Parameters} & \textbf{Returns} & \textbf{Description} \\
  \midrule
  #2
  \bottomrule
  \end{tabular}
  \end{table}
}
\newcommand{\apirow}[4]{%
  \code{#1} & \texttt{#2} & \texttt{#3} & #4 \\
}

% ── Crop marks for print production ─────────────────────────
%  IngramSpark print-ready: crop marks enabled with 0.125" bleed
%  Paper: Cream/natural 60lb  |  Trim: 7.44" × 9.69" (Crown Quarto)
%  Spine width (698 pages, cream 60lb): ~1.12"
\trimmarks  % Adds crop marks at corners
%
%  Spine width reference (698 pages):
%    50lb white: 0.89" (23mm)  |  60lb white: 1.05" (27mm)  |  Cream 60lb: 1.12" (28mm)

% ── PDF metadata ────────────────────────────────────────────
\hypersetup{
  pdftitle={Rust Blockchain: A Full-Stack Implementation Guide},
  pdfauthor={Bill Kunyiha},
  pdfsubject={Blockchain implementation in Rust with Tokio, Axum, Iced, Tauri 2, Docker, and Kubernetes},
  pdfkeywords={Rust, blockchain, Bitcoin, cryptocurrency, distributed systems, Docker, Kubernetes},
  pdfcreator={LaTeX with memoir class},
  pdfproducer={pdfTeX},
}

% ── End of preamble ──────────────────────────────────────────
